\documentclass[12pt, letterpaper]{article}

%% Build from the paper/ directory:
%%   cd paper && pdflatex main.tex
%% or use the root build system:
%%   make paper          (from repo root)
%%   ./build.sh paper    (POSIX)
%%   build.cmd paper     (Windows cmd)

% packages.tex
\usepackage{amsmath, amssymb, amsthm}
\usepackage{graphicx}
\usepackage{xcolor}
\usepackage[
    hidelinks,
    pdftitle={Your Paper Title},
    pdfauthor={Author Name},
    pdfsubject={Paper subject},
    pdfkeywords={keyword1, keyword2, keyword3}
]{hyperref}
\usepackage{booktabs}
\usepackage{array}
\usepackage{geometry}
\usepackage{adjustbox}
\usepackage{longtable}
\usepackage{setspace}
\usepackage{caption}
\usepackage{subcaption}
\usepackage{float}
\onehalfspacing

% commands.tex -- mathematical shorthand for this paper

% Theorem environments (numbered within sections)
\newtheorem{theorem}{Theorem}[section]
\newtheorem{lemma}[theorem]{Lemma}
\newtheorem{corollary}[theorem]{Corollary}
\newtheorem{proposition}[theorem]{Proposition}

% Definitions and remarks (unnumbered or section-numbered as you prefer)
\theoremstyle{definition}
\newtheorem{definition}[theorem]{Definition}
\theoremstyle{remark}
\newtheorem{remark}[theorem]{Remark}

% TODO marker macro -- visible red [TODO: ...] in draft, suppress for release.
% Defined with \long so the argument may contain \par / itemize / multi-
% paragraph prose.  Uses {\color{...} ...} grouping rather than \textcolor
% because the latter is not \long.
\makeatletter
\long\def\todomark#1{{\color{red}[TODO: #1]}}
\long\def\placeholder#1{\par\noindent\todomark{#1}\par}
\makeatother
% To suppress all TODO markers for a release build, uncomment:
% \long\def\todomark#1{}
% \long\def\placeholder#1{}

% Add framework-specific shorthand below.  Keep names short, semantic, and
% consistent across sections.  Examples:
%   \newcommand{\R}{\mathbb{R}}                          % a math symbol
%   \newcommand{\op}[1]{\hat{#1}}                        % a wrapper
%   \newcommand{\eqdef}{\stackrel{\mathrm{def}}{=}}      % a relation


%% Page geometry (also set here as a fallback in case the package path fails)
\geometry{top=1in, bottom=1in, left=1in, right=1in}

%% ---------------------------------------------------------------
%% TITLE PAGE
%% ---------------------------------------------------------------
%% TODO: replace title, subtitle, author, ORCID, email.
%% TODO: at release time, fill in the \thanks{} block with version,
%%       Zenodo DOI, GitHub URL, and any series identifier.

\title{Your Paper Title:\\
       Optional Subtitle on a Second Line%
       \thanks{Version 0.1-DRAFT.
       %% Fill in at release time:
       %% \href{https://doi.org/10.5281/zenodo.XXXXXXXX}%
       %%      {DOI: 10.5281/zenodo.XXXXXXXX}.
       %% \href{https://github.com/USER/REPO}%
       %%      {GitHub: USER/REPO}.
       Pre-release draft; DOI to be assigned.}}
\author{Author Name\\
        \small City, State\\
        \small ORCID: \href{https://orcid.org/0000-0000-0000-0000}{0000-0000-0000-0000}\\
        \small \href{mailto:author@example.com}{author@example.com}}
\date{\today}

\begin{document}

\maketitle

%% ---------------------------------------------------------------
%% FRONT-MATTER NOTICES
%% ---------------------------------------------------------------
\noindent\textbf{License.}  Paper text and figures: Creative Commons
Attribution 4.0 (CC BY 4.0).  Source code in the companion
repository: MIT License.

\medskip
\noindent\textbf{Data availability.}  Numerical experiments, data
files, and the LaTeX source for this paper are publicly available at
\href{https://github.com/USER/REPO}%
     {\nolinkurl{github.com/USER/REPO}}.
%% At release time, point to the tag and the Zenodo DOI:
%% The exact commit corresponding to this release is the vX.Y tag,
%% archived at
%% \href{https://doi.org/10.5281/zenodo.XXXXXXXX}%
%%      {\nolinkurl{doi:10.5281/zenodo.XXXXXXXX}}.

\medskip
\noindent\textbf{Keywords.}  keyword one; keyword two; keyword three.

\newpage

\begin{abstract}
    % abstract.tex
\placeholder{Write the abstract here. Aim for one paragraph: what the paper
claims, what evidence supports it, and what is in scope vs.\ deferred.
End with the falsifiable prediction or the operational test.}

\end{abstract}

%% ---------------------------------------------------------------
%% RELEASE NOTES (only at release time)
%% ---------------------------------------------------------------
%% \bigskip
%% \noindent\textbf{Release notes for vX.Y (since vX.Y-1).}
%% \begin{itemize}
%% \item Item 1.
%% \item Item 2.
%% \end{itemize}

%% \emph{Status legend used throughout.}  \texttt{PASS} / \texttt{PART}
%% / \texttt{STUB} / \texttt{FAIL} -- see Table~\ref{tab:audit}.

\newpage
\tableofcontents
\newpage

%% ============================================================
%% PART I: <name>
%% ============================================================
\section{Introduction}
% introduction.tex
% EXEMPLAR section showing the patterns used throughout the template:
%   - top comment header naming the section and (where applicable) the
%     experiments / derivations it depends on
%   - \label{sec:...} immediately after the section start in main.tex
%   - subsection labels of the form \label{subsec:...}
%   - cross-references via \ref{} / \autoref{}
%   - figure include via \input{...} of a figure-fragment file
%   - bibliography citations with \cite{key}
%   - audit-table reference to Table~\ref{tab:audit}
%
% Replace the prose with your own; keep the labelling and cross-reference
% conventions.

\label{sec:introduction}

\placeholder{One- or two-paragraph orientation: what problem the paper
addresses, why existing approaches fall short, and what the paper's
contribution is. The introduction should be readable without the reader
having internalised any framework-specific notation yet.}

\subsection{Scope}
\label{subsec:scope}

\placeholder{What is in scope for this paper, and what is explicitly
deferred to companion papers or follow-on work. Reference the audit
table (Table~\ref{tab:audit}) if you maintain one --- it is the
canonical answer to ``has this claim been verified yet?''}

%% The audit table is included here so the reader meets it before any
%% claim that depends on a status entry.  Move it elsewhere (e.g. into
%% a "Predictions" section) if that fits the paper better.  Delete this
%% \input if you are not maintaining an audit table.
% ============================================================
% audit_table.tex
%
% A "system audit" table mapping the paper's claims to the evidence
% backing each one.  This is *the* canonical source of truth for which
% claims are PASS / PART / STUB / FAIL --- prose elsewhere in the paper
% should not contradict it.
%
% Status semantics (paste into the abstract or front-matter as well):
%   PASS  -- derivation complete or experiment confirms the claim to
%            stated precision.
%   PART  -- mechanism demonstrated, full quantitative match pending;
%            specific gaps documented in the evidence column.
%   STUB  -- analytical result with no numerical confirmation yet,
%            or planned experiment not yet run.
%   FAIL  -- tested and disconfirmed.
%
% Use \texttt{} for code/experiment names and file paths inside cells.
% Long cells are fine -- the column widths are tuned for prose.
% ============================================================

\label{tab:audit}

{\scriptsize
\setlength{\tabcolsep}{4pt}
\renewcommand{\arraystretch}{0.95}
\begin{longtable}{@{}p{2.6cm}p{3.8cm}p{3.2cm}p{4.4cm}p{1.0cm}@{}}
\caption{\small System audit: paper claims mapped to supporting
derivations and experiments.  Status semantics in the front-matter
(\texttt{PASS} / \texttt{PART} / \texttt{STUB} / \texttt{FAIL}).}
\label{tab:audit_caption} \\
\hline
\textbf{Claim} & \textbf{Mechanism} & \textbf{Continuum / formal limit} & \textbf{Evidence} & \textbf{Status} \\
\hline
\endfirsthead

\multicolumn{5}{l}{\textit{(Table~\ref{tab:audit}, continued from previous page)}} \\
\hline
\textbf{Claim} & \textbf{Mechanism} & \textbf{Continuum / formal limit} & \textbf{Evidence} & \textbf{Status} \\
\hline
\endhead

\hline
\multicolumn{5}{r}{\textit{(continued on next page)}} \\
\endfoot

\hline
\endlastfoot

% ── Example row 1 (replace with a real claim) ───────────────────────
Example claim, fully derived
    & One-line summary of the mechanism
    & The standard / continuum statement it reduces to
    & Derived analytically (\S\ref{subsec:cross_ref_example}, Eq.~\ref{eq:example_identity})
    & \texttt{PASS} \\

% ── Example row 2 (replace with a real claim) ───────────────────────
Example claim, partial evidence
    & One-line summary
    & The target statement
    & \texttt{exp\_XX}: describe the gap concretely (e.g.\ ``mechanism
      reproduced; quantitative match pending higher-resolution run'')
    & \texttt{PART} \\

\hline
\end{longtable}
}


\subsection{A Worked Cross-Reference Example}
\label{subsec:cross_ref_example}

A statement supported by an external citation looks like this~\cite{example2024}.
A statement supported by an in-paper derivation looks like this
(see~\autoref{subsec:scope}).  An equation that is referenced
elsewhere is given a label:
%
\begin{equation}
    \label{eq:example_identity}
    e^{i\pi} + 1 = 0.
\end{equation}
%
Equation~\eqref{eq:example_identity} is referenced by label, not by
hard-coded number.

\subsection{A Worked Figure-Include Example}
\label{subsec:figure_example}

Figures live in \texttt{paper/figures/} (or in the repo-root
\texttt{figures/} directory if the figure is shared with experiments).
The standard pattern is to keep the \verb|\includegraphics| line and the
caption in a small fragment file under \texttt{paper/figures/}, then
\verb|\input| that fragment from inside a \verb|figure| environment
here.  See \texttt{paper/figures/example\_figure.tex} for the template;
once you replace the placeholder with a real figure, uncomment the
include block below.

% \begin{figure}[H]
%     \centering
%     % example_figure.tex
% Figure-fragment file. Convention: \includegraphics on line 1, then the
% caption and label.  Included from a section via:
%   \begin{figure}[H]
%       \centering
%       % example_figure.tex
% Figure-fragment file. Convention: \includegraphics on line 1, then the
% caption and label.  Included from a section via:
%   \begin{figure}[H]
%       \centering
%       \input{figures/example_figure}
%   \end{figure}
% Putting the caption next to the includegraphics keeps the figure file
% self-contained -- if you reuse the same plot in a talk, you take this
% one .tex file with you.

\includegraphics[width=0.7\linewidth]{figures/example_figure.pdf}
\caption[Short list-of-figures caption]{%
    Long caption with full prose explanation of axes, units, and what
    the reader should take away from the figure.  This is the version
    that appears under the figure in the body.}
\label{fig:example}

%   \end{figure}
% Putting the caption next to the includegraphics keeps the figure file
% self-contained -- if you reuse the same plot in a talk, you take this
% one .tex file with you.

\includegraphics[width=0.7\linewidth]{figures/example_figure.pdf}
\caption[Short list-of-figures caption]{%
    Long caption with full prose explanation of axes, units, and what
    the reader should take away from the figure.  This is the version
    that appears under the figure in the body.}
\label{fig:example}

% \end{figure}


\section{A Section Showing the Patterns}
% section_template.tex
% Copy this file to a new name (e.g. sections/dynamics.tex) for each
% additional section, and add a \section{...} \input{sections/...} block
% in main.tex.

\label{sec:template}

\placeholder{Section body. Use \texttt{\textbackslash subsection},
\texttt{\textbackslash subsubsection}, and theorem/lemma environments
as needed.}

\subsection{A Subsection}
\label{subsec:template_a}

\placeholder{Subsection body.}

\begin{theorem}
    \label{thm:template}
    \placeholder{Statement of theorem.}
\end{theorem}

\begin{proof}
    \placeholder{Proof.}
\end{proof}


%% Add more \section{...} \input{sections/...} blocks as needed.
%% Group with comment dividers (e.g., "PART II: DYNAMICS") to keep
%% main.tex skimmable as the paper grows.

\section{Conclusion}
% conclusion.tex

\placeholder{Summarise what was established, what is partial, and what
the natural next paper / experiment would be. Avoid rehashing the
introduction --- the conclusion is for what changed by the end.}


\section*{Acknowledgements}
% acknowledgements.tex

\placeholder{Acknowledge collaborators, reviewers, funding, and tools.
Keep it brief; one paragraph is typical.}


%% ============================================================
%% APPENDICES
%% ============================================================
\appendix

\section{Code and Data}
% code_and_data.tex
\label{sec:code_and_data}

\placeholder{Describe what is in the companion repository: directory
layout, how to inspect the data files, where the figure-generating
scripts live. The aim is that a reader can reproduce or audit any claim
in the paper from this appendix plus the linked DOI.}


\section{Reproducibility Quickstart}
% reproducibility.tex
\label{sec:reproducibility}

\placeholder{Compact reproducibility quickstart. Goal: a reader who has
just cloned the repository at the release tag can reproduce the
paper's headline claim in under five minutes.

Suggested structure:
\begin{itemize}
\item \textbf{Environment.} Operating system, language version, key
      dependencies (point to the manifest file rather than re-listing).
\item \textbf{Build the paper.} The commands needed from a fresh clone.
\item \textbf{Verify the headline result.} A minimal snippet --- ideally
      under ten lines --- that re-derives the central numerical claim.
\item \textbf{Where to look next.} Pointers to the experiment scripts
      and data files for deeper checks.
\end{itemize}}


%% ============================================================
\bibliographystyle{unsrt}
\bibliography{paper-bib/references}

\end{document}
